% Reader question: how can overlapping episodes create continuity without a
% permanent substrate? Claim: direct connectedness is local (adjacent
% incarnations share one witnessed outcome) but continuity is the transitive
% closure of that overlap, not a single unbroken memory. Evidence: three
% incarnations sharing witnesses o_2 (1-2) and o_3 (2-3), so o_1 and o_4 are
% linked without either incarnation directly witnessing the other's outcome
% (Sec. continuity). Must distinguish: a real overlap chain from a repeated
% label with no shared witness. Grammar: an interval/provenance chain,
% rejected: snake/ribbon art (decorative, no evidence marks) and an unlabeled
% chain of circles (no witness identity to check). Five-second test: a reader
% traces o_1 to o_4 through two shared marks and says "linked, but never
% directly."
\begin{figure}[H]
\centering
\pdfigurehue{pdviolet}%
\begin{tikzpicture}[pd figure,x=1cm,y=1cm]
  \node[pd title,anchor=west] at (-1.00,8.80)
    {direct connectedness is local; continuity is the transitive closure};
  % ---- three incarnation intervals, generously separated -------------------
  \draw[pd hairline,line width=5.5pt] (1.40,7.50) -- (4.00,7.50);
  \node[pd label,anchor=east,text width=2.6cm,align=right] at (1.15,7.50) {incarnation 1};
  \draw[pd hairline,line width=5.5pt] (3.40,6.10) -- (8.00,6.10);
  \node[pd label,anchor=east,text width=2.6cm,align=right] at (3.15,6.10) {incarnation 2};
  \draw[pd hairline,line width=5.5pt] (6.60,4.70) -- (10.00,4.70);
  \node[pd label,anchor=east,text width=2.6cm,align=right] at (6.35,4.70) {incarnation 3};
  % ---- witness marks; shared witnesses connect two rows ---------------------
  \node[pd datum] at (2.40,7.50) {};
  \node[pd label,anchor=south] at (2.40,7.78) {$o_1$};
  \node[pd datum] at (4.60,7.50) {};
  \node[pd datum] at (4.60,6.10) {};
  \draw[pd legible rule] (4.60,6.10) -- (4.60,7.50);
  \node[pd label,anchor=south] at (4.60,7.78) {$o_2$};
  \node[pd legible label,anchor=west] at (4.90,6.90) {shared witness $o_2$};
  \node[pd datum] at (7.40,6.10) {};
  \node[pd datum] at (7.40,4.70) {};
  \draw[pd legible rule] (7.40,4.70) -- (7.40,6.10);
  \node[pd label,anchor=west] at (7.68,6.46) {$o_3$};  % see also shared-witness label below, well clear
  \node[pd legible label,anchor=east,text width=2.3cm,align=right] at (7.15,5.55) {shared witness $o_3$};
  \node[pd datum] at (9.20,4.70) {};
  \node[pd label,anchor=south] at (9.20,4.98) {$o_4$};
  % ---- derived accountable chain, drawn once as its own rail ----------------
  \draw[pd legible rule,line width=1.4pt] (2.40,3.40) -- (9.20,3.40);
  \foreach \x in {2.40,4.60,7.40,9.20} \node[pd legible datum] at (\x,3.40) {};
  \node[pd legible label,anchor=north,text width=8.0cm,align=center] at (5.80,3.10)
    {transitive ledger rail: $o_1\leadsto o_2\leadsto o_3\leadsto o_4$};
  % ---- counterexample: a repeated name without an overlapping witness -------
  \draw[pd breach rule,dashed] (1.40,1.60) -- (4.00,1.60);
  \draw[pd breach rule,dashed] (4.80,1.60) -- (7.60,1.60);
  \draw[pd breach rule] (4.25,1.45) -- (4.55,1.75);
  \draw[pd breach rule] (4.25,1.75) -- (4.55,1.45);
  \node[pd breach label,anchor=north,text width=8.0cm,align=center] at (4.50,1.15)
    {same label, no shared witness $\Rightarrow$ no inherited continuity};
\end{tikzpicture}
\caption{\textbf{Parfitian continuity is an overlap chain, not permanent
connectedness.} Incarnation 1 shares witnessed outcome $o_2$ with incarnation 2;
incarnation 2 shares $o_3$ with incarnation 3. No later body must directly
remember $o_1$: the overlapping evidence establishes the transitive ledger rail
$o_1\leadsto o_2\leadsto o_3\leadsto o_4$. The dashed counterexample repeats a
label without a shared witness and therefore inherits nothing. [internal]}
\label{fig:parfit}
\end{figure}
